% =========================================================================
%  Figura (Problema 3.3): los dos estados del camino A con a = 1, b = 2.
%  phi1 = 1, phi2 = 0: phi = 1 (r <= 1), 2/r - 1 (1..2), 0 (r >= 2)
%  phi1 = 0, phi2 = 1: phi = 0 (r <= 1), 2 - 2/r (1..2), 2/r (r >= 2)
% =========================================================================
\begin{figure}[t]
\centering
\begin{tikzpicture}[font=\small]
  \begin{axis}[tema, name=a, width=7.0cm, height=4.8cm, xmin=0, xmax=4.2, ymin=0, ymax=1.18,
      xlabel={$r/a$}, ylabel={$\phi$}, xtick={0,1,2,3,4}, ytick={0,0.5,1},
      title={\footnotesize $\phi_1=1,\ \phi_2=0$: columna 1 de $\Cm$}]
    \fill[azulvelo!60] (axis cs:0,0) rectangle (axis cs:1,1.18);
    \addplot[azulprof, line width=1.2pt] coordinates {(0,1) (1,1)};
    \addplot[azulprof, line width=1.2pt, domain=1:2, samples=40] {2/x-1};
    \addplot[azulprof, line width=1.2pt] coordinates {(2,0) (4.2,0)};
    \draw[punteado] (axis cs:2,0) -- (axis cs:2,1.18);
    \node[font=\scriptsize, anchor=west, align=left] at (axis cs:2.2,0.65) {$Q_1=c_{11}>0$\\ $Q_2=c_{21}=-c_{11}$};
  \end{axis}
  \node[anchor=north] at ($(a.outer south)+(0,-0.05)$) {(a)};
  \begin{axis}[tema, name=b, at={($(a.east)+(1.4cm,0)$)}, anchor=west, width=7.0cm, height=4.8cm,
      xmin=0, xmax=4.2, ymin=0, ymax=1.18,
      xlabel={$r/a$}, ylabel={$\phi$}, xtick={0,1,2,3,4}, ytick={0,0.5,1},
      title={\footnotesize $\phi_1=0,\ \phi_2=1$: columna 2 de $\Cm$}]
    \fill[azulvelo!60] (axis cs:0,0) rectangle (axis cs:1,1.18);
    \addplot[azulacero, line width=1.2pt] coordinates {(0,0) (1,0)};
    \addplot[azulacero, line width=1.2pt, domain=1:2, samples=40] {2-2/x};
    \addplot[azulacero, line width=1.2pt, domain=2:4.2, samples=40] {2/x};
    \draw[punteado] (axis cs:2,0) -- (axis cs:2,1.18);
    \node[font=\scriptsize, anchor=west, align=left] at (axis cs:2.45,0.38) {$Q_1=c_{12}<0$\\ $Q_2=c_{22}>0$};
  \end{axis}
  \node[anchor=north] at ($(b.outer south)+(0,-0.05)$) {(b)};
\end{tikzpicture}
\caption{Problema 3.3, camino A ($b=2a$). En cada estado se fija el potencial de un conductor a la unidad y el otro a tierra; las cargas que aparecen son una
columna de la matriz de coeficientes de capacidad e influencia. (a) La esfera interior a potencial unidad: todo su flujo acaba en la exterior. (b) La esfera
exterior a potencial unidad: parte de su flujo va al infinito y parte induce carga negativa en la interior, a tierra.}
\label{fig:sem-esferico}
\end{figure}
